% !TEX program = pdflatex
\documentclass[10pt]{article}

\usepackage[utf8]{inputenc}
\usepackage[T1]{fontenc}
\usepackage[english]{babel}
\usepackage{lmodern}
\usepackage{microtype}
\usepackage[margin=0.72in]{geometry}
\usepackage{amsmath,amssymb,amsthm,mathtools}
\usepackage{booktabs}
\usepackage{xcolor}
\usepackage{enumitem}
\usepackage[most]{tcolorbox}
\usepackage{fancyhdr}
\usepackage[hidelinks]{hyperref}

\definecolor{CourseBlue}{HTML}{1A3C68}
\definecolor{CourseGreen}{HTML}{216E39}
\definecolor{CourseGold}{HTML}{B8860B}
\definecolor{CourseRed}{HTML}{9A2C2C}
\definecolor{SoftBlue}{HTML}{F2F6FA}
\definecolor{SoftGreen}{HTML}{F1F8F3}
\definecolor{SoftGold}{HTML}{FFF8E8}
\definecolor{SoftRed}{HTML}{FFF3F3}

\setlength{\parindent}{0pt}
\setlength{\parskip}{0.32em}
\setlist[itemize]{leftmargin=1.5em,itemsep=0.12em,topsep=0.2em}
\setlist[enumerate]{leftmargin=1.7em,itemsep=0.12em,topsep=0.2em}
\allowdisplaybreaks

\pagestyle{fancy}
\fancyhf{}
\lhead{Stochastic Calculus and Option Pricing}
\rhead{Lecture 5 Notes}
\cfoot{\thepage}
\renewcommand{\headrulewidth}{0.3pt}

\newcommand{\E}{\mathbb E}
\newcommand{\Pp}{\mathbb P}
\newcommand{\R}{\mathbb R}
\newcommand{\1}{\mathbf 1}
\newcommand{\cF}{\mathcal F}
\newcommand{\dd}{\,\mathrm d}
\newcommand{\Ito}{\mathrm{It\hat{o}}}
\newcommand{\Ncdf}{\mathrm N}
\newcommand{\phiN}{\varphi}
\DeclareMathOperator{\Var}{Var}
\DeclareMathOperator{\Cov}{Cov}

\theoremstyle{definition}
\newtheorem{definition}{Definition}[section]
\newtheorem{proposition}[definition]{Proposition}
\newtheorem{theorem}[definition]{Theorem}
\newtheorem{lemma}[definition]{Lemma}
\newtheorem{corollary}[definition]{Corollary}
\newtheorem{remark}[definition]{Remark}

\newtcolorbox{boardexample}[1]{
  enhanced,breakable,colback=SoftBlue,colframe=CourseBlue,
  boxrule=0.55pt,arc=1mm,left=1.5mm,right=1.5mm,top=1mm,bottom=1mm,
  before skip=0.55em,after skip=0.55em,
  title={Example --- #1},fonttitle=\bfseries}
\newtcolorbox{studentexercise}[1]{
  enhanced,breakable,colback=SoftGold,colframe=CourseGold,
  boxrule=0.55pt,arc=1mm,left=1.5mm,right=1.5mm,top=1mm,bottom=1mm,
  before skip=0.55em,after skip=0.35em,
  title={Exercise --- #1},fonttitle=\bfseries}
\newtcolorbox{shortsolution}[1][]{
  enhanced,breakable,colback=SoftGreen,colframe=CourseGreen,
  boxrule=0.55pt,arc=1mm,left=1.5mm,right=1.5mm,top=1mm,bottom=1mm,
  before skip=0.25em,after skip=0.6em,
  title={Short solution#1},fonttitle=\bfseries}
\newtcolorbox{keyidea}[1]{
  enhanced,breakable,colback=white,colframe=CourseGreen,
  boxrule=0.7pt,arc=1mm,left=1.5mm,right=1.5mm,top=1mm,bottom=1mm,
  before skip=0.5em,after skip=0.5em,
  title={#1},fonttitle=\bfseries}
\newtcolorbox{warningbox}[1]{
  enhanced,breakable,colback=SoftRed,colframe=CourseRed,
  boxrule=0.65pt,arc=1mm,left=1.5mm,right=1.5mm,top=1mm,bottom=1mm,
  before skip=0.5em,after skip=0.5em,
  title={#1},fonttitle=\bfseries}
\newtcolorbox{assumptionbox}[1]{
  enhanced,breakable,colback=SoftGold,colframe=CourseGold,
  boxrule=0.65pt,arc=1mm,left=1.5mm,right=1.5mm,top=1mm,bottom=1mm,
  before skip=0.5em,after skip=0.5em,
  title={#1},fonttitle=\bfseries}

\title{\textbf{Black--Scholes from Replication}\\
\large Lecture 5 Notes --- Self-Financing, PDE, Verification, and Greeks}
\author{}
\date{}

\begin{document}
\maketitle
\vspace{-1.4em}

\begin{keyidea}{Goals and logical structure}
This lecture connects stochastic calculus to option valuation through one
reproducible chain:
\[
\text{\(\Ito\) dynamics}
\longrightarrow
\text{self-financing replication}
\longrightarrow
\text{Black--Scholes PDE}
\longrightarrow
\text{Gaussian expectation and Greeks}.
\]
We first practice \(\Ito\)'s formula on examples not used in Lecture 4.
We then define self-financing carefully in discrete and continuous time,
derive the PDE for a European payoff \(g(S_T)\), solve and verify it, compute
the call and put formulas and their Greeks, and conclude with the practical
limits of continuous-time hedging.
\end{keyidea}

% ============================================================
\section{Three New Uses of the It\^o Formula}
% ============================================================

\subsection{The generator organizes the calculation}

Let
\[
\dd X_t=b(t,X_t)\dd t+a(t,X_t)\dd W_t
\]
and \(f\in C^{1,2}\). The one-dimensional \(\Ito\) formula is
\[
\dd f(t,X_t)
=\left(f_t+bf_x+\frac12a^2f_{xx}\right)(t,X_t)\dd t
+a(t,X_t)f_x(t,X_t)\dd W_t.
\]
It is convenient to define the generator
\[
\mathcal Lf(t,x)
=b(t,x)f_x(t,x)+\frac12a(t,x)^2f_{xx}(t,x),
\]
so that
\[
\boxed{
\dd f(t,X_t)
=\bigl(f_t+\mathcal Lf\bigr)(t,X_t)\dd t
+a(t,X_t)f_x(t,X_t)\dd W_t.}
\]
The operator \(f_t+\mathcal Lf\) is exactly the drift of the transformed
process.

\begin{warningbox}{Driftless is not automatically a true martingale}
If \(f_t+\mathcal Lf=0\), then \(f(t,X_t)\) is first a local martingale.
An integrability, boundedness, or localization argument is needed before
taking expectations or conditional expectations.
\end{warningbox}

\subsection{Dimension creates drift}

\begin{boardexample}{the squared norm of Brownian motion}
Let \(W=(W^1,\ldots,W^d)\) be a standard \(d\)-dimensional Brownian
motion and set
\[
R_t=\|W_t\|^2=\sum_{i=1}^d(W_t^i)^2.
\]
Applying the multidimensional formula to \(f(x)=\|x\|^2\) gives
\[
\boxed{
\dd R_t=2W_t^\top\dd W_t+d\,\dd t.}
\]
Consequently,
\[
R_t-dt=2\int_0^tW_s^\top\dd W_s
\]
is a square-integrable martingale on bounded horizons, and
\[
\E[\|W_t\|^2]=dt.
\]
Each independent Brownian direction contributes one unit of quadratic
variation per unit of time.
\end{boardexample}

\begin{studentexercise}{replace the identity covariance}
Suppose \(X_t=CW_t\), where \(C\) is a deterministic \(m\times d\) matrix.
Find \(\dd\|X_t\|^2\) and \(\E\|X_t\|^2\).
\end{studentexercise}
\begin{shortsolution}
Since \(\dd X_t=C\dd W_t\),
\[
\dd[X^i,X^j]_t=(CC^\top)_{ij}\dd t.
\]
Therefore
\[
\dd\|X_t\|^2
=2X_t^\top C\dd W_t+\operatorname{tr}(CC^\top)\dd t,
\]
and
\[
\boxed{\E\|X_t\|^2=t\,\operatorname{tr}(CC^\top).}
\]
\end{shortsolution}

\subsection{Adding a path variable restores the Markov state}

Define integrated Brownian motion
\[
A_t=\int_0^tW_s\dd s.
\]
The pair \((W_t,A_t)\) is a two-dimensional state process:
\[
\dd W_t=\dd W_t,\qquad \dd A_t=W_t\dd t.
\]
Although \(A_t\) depends on the entire past of \(W\), augmenting the state by
\(A_t\) makes the dynamics local again.

\begin{boardexample}{a product with integrated Brownian motion}
For \(Y_t=W_tA_t\), the product rule gives
\[
\dd Y_t
=A_t\dd W_t+W_t\dd A_t+\dd[W,A]_t.
\]
The process \(A\) has finite variation, so \([W,A]=0\), and hence
\[
\boxed{\dd(W_tA_t)=A_t\dd W_t+W_t^2\dd t.}
\]
Taking expectations,
\[
\boxed{
\E[W_tA_t]
=\int_0^t\E[W_s^2]\dd s
=\frac{t^2}{2}.}
\]
This agrees with the direct covariance calculation
\[
\Cov\left(W_t,\int_0^tW_s\dd s\right)
=\int_0^t\min(t,s)\dd s=\frac{t^2}{2}.
\]
\end{boardexample}

\subsection{A future event becomes a backward-PDE martingale}

Let \(\Ncdf\) and \(\phiN\) denote the standard normal cdf and density. Fix
\(T>0\) and define, for \(t<T\),
\[
u(t,x)=\Ncdf\left(\frac{x}{\sqrt{T-t}}\right).
\]
Writing \(\tau=T-t\) and \(z=x/\sqrt{\tau}\),
\[
u_x=\frac{\phiN(z)}{\sqrt{\tau}},
\qquad
u_{xx}=-\frac{x\phiN(z)}{\tau^{3/2}},
\qquad
u_t=\frac{x\phiN(z)}{2\tau^{3/2}}.
\]
Therefore
\[
u_t+\frac12u_{xx}=0
\]
and
\[
\dd u(t,W_t)
=\frac{\phiN(W_t/\sqrt{T-t})}{\sqrt{T-t}}\dd W_t.
\]
The probabilistic meaning is
\[
\boxed{
u(t,W_t)=\Pp(W_T>0\mid\cF_t).}
\]
Indeed, conditional on \(\cF_t\),
\[
W_T=W_t+\sqrt{T-t}\,Z,\qquad Z\sim N(0,1).
\]

\begin{keyidea}{Preview of Black--Scholes}
A terminal payoff, a backward PDE, a conditional expectation, and a
replicating stochastic integral will again describe the same process.
\end{keyidea}

% ============================================================
\section{Securities, Trading Strategies, and Self-Financing}
% ============================================================

\subsection{Model assumptions and market assumptions}

On a filtered probability space
\((\Omega,\cF,(\cF_t)_{0\leq t\leq T},\Pp)\), consider one risky asset
and one bank account:
\[
\frac{\dd S_t}{S_t}=\mu\dd t+\sigma\dd W_t,
\qquad
\dd B_t=rB_t\dd t,\qquad B_0=1.
\]
Thus \(B_t=e^{rt}\). Throughout the basic Black--Scholes model,
\(\sigma>0\).

\begin{assumptionbox}{Model assumptions}
\begin{itemize}
  \item The stock has continuous geometric Brownian dynamics.
  \item The parameters \(\mu,\sigma,r\) are constant.
  \item The stock pays no dividends and does not default.
  \item One Brownian motion generates all uncertainty in the model.
\end{itemize}
\end{assumptionbox}

\begin{assumptionbox}{Market assumptions}
\begin{itemize}
  \item Trading is continuous and frictionless.
  \item There are no transaction costs, bid--ask spreads, or taxes.
  \item Arbitrary long and short positions are allowed.
  \item Cash can be borrowed or lent at the same rate \(r\).
  \item Strategies are predictable and admissible.
  \item The market satisfies no arbitrage.
\end{itemize}
\end{assumptionbox}

\begin{warningbox}{Keep the logical roles separate}
GBM is a model assumption. No arbitrage is a market principle. Completeness
comes from having enough traded assets to span the modeled shocks. Constant
coefficients and zero dividends are simplifications that produce the familiar
closed form. None of these statements implies all the others.
\end{warningbox}

\subsection{Portfolio value and predictability}

\begin{definition}[Trading strategy]
A strategy consists of predictable processes
\((\Delta_t,\beta_t)\), where \(\Delta_t\) is the number of stock shares and
\(\beta_t\) is the number of bank-account units held at time \(t\). Its value is
\[
\boxed{X_t=\Delta_tS_t+\beta_tB_t.}
\]
\end{definition}

Predictability expresses causality: the position used over the next instant
must be selected from information already available. Suitable admissibility
conditions also exclude doubling strategies and ensure that the stochastic
integrals are well defined.

\subsection{Self-financing in discrete time}

The discrete formulation makes the cash-flow accounting transparent. Let
\(0=t_0<\cdots<t_N=T\). The holdings
\((\Delta_n,\beta_n)\), chosen at time \(t_n\), are kept over
\((t_n,t_{n+1}]\).

Immediately before rebalancing at \(t_n\), the old positions are worth
\[
X_n^-=\Delta_{n-1}S_{t_n}+\beta_{n-1}B_{t_n}.
\]
Immediately after rebalancing, the new positions are worth
\[
X_n=\Delta_nS_{t_n}+\beta_nB_{t_n}.
\]

\begin{definition}[Discrete self-financing condition]
No cash is added or withdrawn at \(t_n\):
\[
X_n=X_n^-.
\]
Equivalently,
\[
\boxed{
S_{t_n}(\Delta_n-\Delta_{n-1})
+B_{t_n}(\beta_n-\beta_{n-1})=0.}
\]
\end{definition}

Buying stock is financed by reducing the cash position; selling stock
generates cash. Rebalancing changes composition, not wealth.

Between \(t_n\) and \(t_{n+1}\), the assets move while the holdings remain
fixed. Hence
\[
\boxed{
X_{n+1}^- -X_n
=\Delta_n(S_{t_{n+1}}-S_{t_n})
+\beta_n(B_{t_{n+1}}-B_{t_n}).}
\]
Portfolio gains come from asset-price changes, not from the act of
rebalancing.

\subsection{Continuous-time self-financing}

\begin{definition}[Continuous self-financing condition]
A strategy is self-financing when
\[
\boxed{
X_t=X_0+\int_0^t\Delta_u\dd S_u+\int_0^t\beta_u\dd B_u.}
\]
In differential shorthand,
\[
\boxed{\dd X_t=\Delta_t\dd S_t+\beta_t\dd B_t.}
\]
\end{definition}

The stochastic integrals are limits of predictable left-point trading gains.
Using
\[
\beta_tB_t=X_t-\Delta_tS_t,
\qquad
\dd B_t=rB_t\dd t,
\]
the gain equation becomes
\[
\boxed{
\dd X_t
=\Delta_t\dd S_t+r(X_t-\Delta_tS_t)\dd t.}
\]

\begin{proposition}[Equivalent rebalancing constraint]
If \(\Delta\) is a continuous semimartingale and
\(X_t=\Delta_tS_t+\beta_tB_t\), then self-financing is equivalent to
\[
\boxed{
S_t\dd\Delta_t+\dd[\Delta,S]_t+B_t\dd\beta_t=0.}
\]
\end{proposition}
\begin{proof}
The product rule gives
\[
\dd X_t
=\Delta_t\dd S_t+S_t\dd\Delta_t+\dd[\Delta,S]_t
+\beta_t\dd B_t+B_t\dd\beta_t.
\]
Comparing this with the self-financing gain equation yields the claim.
\end{proof}

\begin{warningbox}{Value identity versus gain identity}
The identity \(X=\Delta S+\beta B\) says what the portfolio is worth.
The identity
\(\dd X=\Delta\dd S+\beta\dd B\) says how its value changes without external
cash. A dynamic hedge needs both statements.
\end{warningbox}

% ============================================================
\section{Replication and the Black--Scholes PDE}
% ============================================================

\subsection{The European claim and its It\^o dynamics}

Fix a European payoff
\[
H=g(S_T).
\]
Assume that its value can be written as a sufficiently smooth Markov function
\[
V_t=p(t,S_t),
\qquad
p(T,s)=g(s).
\]
The condition \(p\in C^{1,2}([0,T)\times(0,\infty))\) will be used before
maturity. The payoff itself may have a kink at \(T\).

By \(\Ito\)'s formula,
\[
\boxed{
\begin{aligned}
\dd V_t
={}&\left(
p_t+\mu S_tp_s+\frac12\sigma^2S_t^2p_{ss}
\right)(t,S_t)\dd t\\
&+\sigma S_tp_s(t,S_t)\dd W_t.
\end{aligned}}
\]

\subsection{Match a self-financing replicating portfolio}

A replicating strategy satisfies
\[
X_t=p(t,S_t)
\quad\text{for all }t,
\qquad
X_T=g(S_T).
\]
Its self-financing dynamics are
\[
\begin{aligned}
\dd X_t
&=\Delta_t\dd S_t+r(X_t-\Delta_tS_t)\dd t\\
&=\left[rX_t+(\mu-r)S_t\Delta_t\right]\dd t
+\sigma S_t\Delta_t\dd W_t.
\end{aligned}
\]
Matching the Brownian coefficients in \(\dd X_t=\dd V_t\) gives
\[
\boxed{\Delta_t=p_s(t,S_t).}
\]
The stock holding in the replicating strategy is the option delta.

Now match the drift coefficients with \(X_t=p(t,S_t)\):
\[
rp+(\mu-r)sp_s
=p_t+\mu sp_s+\frac12\sigma^2s^2p_{ss}.
\]
The physical drift cancels. We obtain the pricing equation.

\begin{theorem}[Black--Scholes PDE for a European payoff]
Under the preceding model, market, regularity, and replication assumptions,
the value \(p\) satisfies
\[
\boxed{
\begin{cases}
p_t(t,s)+rs\,p_s(t,s)
+\dfrac12\sigma^2s^2p_{ss}(t,s)-rp(t,s)=0,
&t<T,\ s>0,\\[0.4em]
p(T,s)=g(s).
\end{cases}}
\]
The replicating holdings are
\[
\boxed{
\Delta_t=p_s(t,S_t),
\qquad
\beta_t=\frac{p(t,S_t)-S_tp_s(t,S_t)}{B_t}.}
\]
\end{theorem}

\begin{keyidea}{Why \(\mu\) disappears}
The option and the replicating portfolio have the same exposure to the
stock. Their physical expected-return terms are therefore identical and
cancel. Valuation depends on the cost of financing \(r\) and the unhedgeable
quadratic-variation scale \(\sigma\), not on an estimate of \(\mu\).
\end{keyidea}

\begin{warningbox}{What the cancellation does not mean}
The physical drift remains relevant for forecasting, portfolio choice, risk
management, and the empirical distribution of hedging errors. It disappears
from the replication price inside the idealized complete market.
\end{warningbox}

\subsection{No arbitrage and completeness}

The replication argument uses two distinct economic ideas:

\begin{itemize}
  \item \textbf{Completeness:} the one traded stock spans the one Brownian
        source of risk, so the option's diffusion coefficient can be matched.
  \item \textbf{No arbitrage:} a claim and a self-financing portfolio with the
        same terminal payoff must have the same initial price.
\end{itemize}

With jumps or an additional stochastic-volatility shock, one stock generally
cannot match every source of uncertainty. The market may become incomplete,
and no-arbitrage alone may determine a range of prices rather than a unique
one.

\begin{studentexercise}{add a continuous dividend yield}
Suppose the stock pays a continuous dividend yield \(q\). A stock holder
receives \(qS_t\dd t\) in addition to the price gain \(\dd S_t\). Derive the
PDE for a European claim.
\end{studentexercise}
\begin{shortsolution}
The stock position earns
\[
\Delta_t(\dd S_t+qS_t\dd t).
\]
Repeating the self-financing replication argument gives
\[
\boxed{
p_t+(r-q)sp_s+\frac12\sigma^2s^2p_{ss}-rp=0,
\qquad p(T,s)=g(s).}
\]
The no-dividend model is the special case \(q=0\).
\end{shortsolution}

% ============================================================
\section{Solving the Black--Scholes PDE}
% ============================================================

\subsection{Step 1: time to maturity}

We solve
\[
p_t+rsp_s+\frac12\sigma^2s^2p_{ss}-rp=0,
\qquad p(T,s)=g(s).
\]
Let
\[
\tau=T-t,
\qquad
v(\tau,s)=p(T-\tau,s).
\]
Since \(p_t=-v_\tau\),
\[
\boxed{
v_\tau
=rsv_s+\frac12\sigma^2s^2v_{ss}-rv,
\qquad
v(0,s)=g(s).}
\]

\subsection{Step 2: remove the discounting term}

Define
\[
u(\tau,s)=e^{r\tau}v(\tau,s).
\]
Then
\[
u_\tau=ru+e^{r\tau}v_\tau,\qquad
u_s=e^{r\tau}v_s,\qquad
u_{ss}=e^{r\tau}v_{ss}.
\]
The two \(r u\) terms cancel:
\[
\boxed{
u_\tau
=rsu_s+\frac12\sigma^2s^2u_{ss},
\qquad
u(0,s)=g(s).}
\]

\subsection{Step 3: logarithmic coordinates}

Set
\[
x=\log s,
\qquad
w(\tau,x)=u(\tau,e^x).
\]
The derivatives satisfy
\[
w_x=su_s,\qquad
w_{xx}=su_s+s^2u_{ss},
\]
and therefore
\[
u_s=\frac{w_x}{s},
\qquad
u_{ss}=\frac{w_{xx}-w_x}{s^2}.
\]
Substitution yields
\[
\boxed{
w_\tau
=\frac12\sigma^2w_{xx}
+\left(r-\frac12\sigma^2\right)w_x,
\qquad
w(0,x)=g(e^x).}
\]
The coefficient \(r-\sigma^2/2\) is the drift of the log price under the
pricing dynamics.

\subsection{Step 4: remove transport and solve the heat equation}

Let
\[
a=r-\frac12\sigma^2,
\qquad
\varphi(\tau,y)=w(\tau,y-a\tau).
\]
Then
\[
\varphi_\tau=w_\tau-aw_x,\qquad
\varphi_y=w_x,\qquad
\varphi_{yy}=w_{xx},
\]
so
\[
\boxed{
\varphi_\tau=\frac12\sigma^2\varphi_{yy},
\qquad
\varphi(0,y)=g(e^y).}
\]
If \(Z\sim N(0,1)\), the heat-kernel solution is
\[
\varphi(\tau,y)
=\E\left[g\left(e^{y+\sigma\sqrt{\tau}Z}\right)\right].
\]
Equivalently,
\[
\varphi(\tau,y)
=\int_{\R}g(e^z)
\frac{\exp\left(-\frac{(z-y)^2}{2\sigma^2\tau}\right)}
{\sqrt{2\pi\sigma^2\tau}}\dd z.
\]

\begin{theorem}[General European-payoff formula]
Let \(\tau=T-t\). Under suitable growth and regularity assumptions,
\[
\boxed{
\begin{aligned}
p(t,s)
=e^{-r\tau}\E\Bigg[
g\Bigg(
s\exp\Big\{&
\left(r-\frac12\sigma^2\right)\tau\\
&+\sigma\sqrt{\tau}Z
\Big\}
\Bigg)
\Bigg],
\qquad Z\sim N(0,1).
\end{aligned}}
\]
\end{theorem}

\begin{remark}
This is a general formula, but not necessarily an elementary closed form.
For an arbitrary payoff \(g\), it reduces valuation to a one-dimensional
Gaussian integral.
\end{remark}

\begin{warningbox}{Regularity at maturity}
Calls, puts, and digital payoffs are not \(C^2\) at maturity. Gaussian
convolution smooths their prices for every \(t<T\). A fully general treatment
distinguishes classical, weak, and viscosity solutions; here we use the
classical formula before maturity and recover the payoff in the terminal
limit.
\end{warningbox}

\begin{studentexercise}{price a power payoff}
For \(g(s)=s^m\), with \(m\in\R\) such that the expectation is finite,
compute \(p(t,s)\).
\end{studentexercise}
\begin{shortsolution}
Using \(\E[e^{cZ}]=e^{c^2/2}\),
\[
\begin{aligned}
p(t,s)
&=e^{-r\tau}s^m
\exp\left(
m(r-\tfrac12\sigma^2)\tau
+\tfrac12m^2\sigma^2\tau
\right)\\
&=\boxed{
s^m\exp\left(
(m-1)r\tau+\frac12m(m-1)\sigma^2\tau
\right).}
\end{aligned}
\]
The cases \(m=0\) and \(m=1\) give a zero-coupon payment and one share,
respectively.
\end{shortsolution}

% ============================================================
\section{Verification: PDE, Expectation, and Hedge}
% ============================================================

\subsection{The discounted candidate must be a martingale}

Introduce an auxiliary stock, started at \(S_t^r=s\), with dynamics
\[
\frac{\dd S_u^r}{S_u^r}
=r\dd u+\sigma\dd W_u,
\qquad u\in[t,T].
\]
Its terminal value is
\[
S_T^r
=s\exp\left(
\left(r-\frac12\sigma^2\right)(T-t)
+\sigma(W_T-W_t)
\right).
\]

\begin{theorem}[Verification theorem]
Suppose \(p\in C^{1,2}\) solves the Black--Scholes PDE with terminal
condition \(g\), and assume the stochastic integral below is a true
martingale. Then
\[
\boxed{
p(t,s)
=\E\left[
e^{-r(T-t)}g(S_T^r)\mid S_t^r=s
\right].}
\]
Moreover, in the original market, the strategy
\[
\Delta_t=p_s(t,S_t),
\qquad
\beta_t=\frac{p(t,S_t)-S_tp_s(t,S_t)}{B_t}
\]
is self-financing and replicates \(g(S_T)\).
\end{theorem}

\begin{proof}
Apply \(\Ito\)'s formula and the product rule:
\[
\begin{aligned}
\dd\left(e^{-r(u-t)}p(u,S_u^r)\right)
={}&e^{-r(u-t)}
\left[
p_t+rS_u^rp_s
+\frac12\sigma^2(S_u^r)^2p_{ss}-rp
\right]\dd u\\
&+e^{-r(u-t)}\sigma S_u^rp_s\dd W_u.
\end{aligned}
\]
The drift vanishes by the PDE. Taking conditional expectations between
\(t\) and \(T\) and using \(p(T,\cdot)=g\) proves the representation.

For replication under the original dynamics, rearrange the PDE inside
\(\Ito\)'s formula:
\[
\dd p(t,S_t)
=p_s(t,S_t)\dd S_t
+r\bigl(p(t,S_t)-S_tp_s(t,S_t)\bigr)\dd t.
\]
This equals
\(\Delta_t\dd S_t+\beta_t\dd B_t\), and the terminal value equals
\(g(S_T)\).
\end{proof}

\begin{keyidea}{The logical loop}
\[
\boxed{
\text{self-financing replication}
\Longleftrightarrow
\text{Black--Scholes PDE}
\Longleftrightarrow
\text{discounted expectation}.}
\]
Replication explains uniqueness and identifies the hedge. The expectation
computes the PDE solution. Verification checks that both descriptions agree.
\end{keyidea}

% ============================================================
\section{The European Call and Put}
% ============================================================

\subsection{A truncated Gaussian identity}

Let
\[
\phiN(z)=\frac{e^{-z^2/2}}{\sqrt{2\pi}},
\qquad
\Ncdf(x)=\int_{-\infty}^x\phiN(z)\dd z.
\]

\begin{lemma}[Exponential tilting of a Gaussian tail]
For \(a,b\in\R\) and \(Z\sim N(0,1)\),
\[
\boxed{
\E\left[e^{aZ}\1_{\{Z>b\}}\right]
=e^{a^2/2}\Ncdf(a-b).}
\]
\end{lemma}
\begin{proof}
Complete the square:
\[
e^{az}\phiN(z)
=e^{a^2/2}\phiN(z-a).
\]
Integrating over \((b,\infty)\) and substituting \(y=z-a\) gives the result.
\end{proof}

\subsection{The exercise event produces \texorpdfstring{\(d_-\)}{d-minus}}

For a call,
\[
g(s)=(s-K)^+,
\qquad
\tau=T-t.
\]
Under the auxiliary drift-\(r\) dynamics,
\[
S_T^r
=s\exp\left(
\left(r-\frac12\sigma^2\right)\tau
+\sigma\sqrt{\tau}Z
\right).
\]
The price splits as
\[
C(t,s)
=e^{-r\tau}
\left(
\E[S_T^r\1_{\{S_T^r>K\}}]
-K\Pp(S_T^r>K)
\right).
\]
Now
\[
S_T^r>K
\iff
Z>-d_-,
\]
where
\[
\boxed{
d_-=
\frac{\log(s/K)+(r-\frac12\sigma^2)\tau}
{\sigma\sqrt{\tau}}.}
\]
Therefore
\[
\Pp(S_T^r>K)=\Ncdf(d_-).
\]

\subsection{The stock-weighted term produces \texorpdfstring{\(d_+\)}{d-plus}}

Apply the truncated Gaussian lemma with
\(a=\sigma\sqrt{\tau}\) and \(b=-d_-\):
\[
\E\left[
e^{\sigma\sqrt{\tau}Z}\1_{\{Z>-d_-\}}
\right]
=e^{\sigma^2\tau/2}
\Ncdf(d_-+\sigma\sqrt{\tau}).
\]
Define
\[
\boxed{
d_+=d_-+\sigma\sqrt{\tau}
=\frac{\log(s/K)+(r+\frac12\sigma^2)\tau}
{\sigma\sqrt{\tau}}.}
\]
It follows that
\[
\E[S_T^r\1_{\{S_T^r>K\}}]
=se^{r\tau}\Ncdf(d_+).
\]

\begin{theorem}[Black--Scholes call and put formulas]
For \(\tau=T-t\),
\[
\boxed{
C(t,s)
=s\Ncdf(d_+)-Ke^{-r\tau}\Ncdf(d_-),}
\]
where
\[
\boxed{
d_\pm
=\frac{\log(s/K)+r\tau}{\sigma\sqrt{\tau}}
\pm\frac12\sigma\sqrt{\tau}.}
\]
The put price is
\[
\boxed{
P(t,s)
=Ke^{-r\tau}\Ncdf(-d_-)-s\Ncdf(-d_+).}
\]
\end{theorem}

\begin{proof}[Put formula]
The terminal identity
\[
(S_T-K)^+-(K-S_T)^+=S_T-K
\]
is replicated by one share minus a zero-coupon payment \(K\). No arbitrage
gives put--call parity:
\[
C(t,s)-P(t,s)=s-Ke^{-r\tau}.
\]
Insert the call formula and use
\(\Ncdf(-x)=1-\Ncdf(x)\).
\end{proof}

\begin{warningbox}{Do not interpret both normal terms as the same probability}
\(\Ncdf(d_-)\) is the exercise probability under the drift-\(r\) pricing
dynamics. The factor \(\Ncdf(d_+)\) comes from a stock-weighted truncated
expectation. Its later interpretation requires a stock-numeraire measure.
\end{warningbox}

\begin{boardexample}{boundary and maturity checks}
The formula satisfies
\[
\begin{array}{lll}
s\downarrow0: & d_\pm\to-\infty, & C(t,s)\to0,\\
s\uparrow\infty: & d_\pm\to+\infty, & C(t,s)\sim s-Ke^{-r\tau},\\
\tau\downarrow0: & & C(t,s)\to(s-K)^+.
\end{array}
\]
It also satisfies the no-arbitrage bounds
\[
(s-Ke^{-r\tau})^+\leq C(t,s)\leq s.
\]
These checks catch sign errors in \(r\tau\), discounting, and \(d_\pm\).
\end{boardexample}

\begin{studentexercise}{a cash-or-nothing digital}
Find the price of the payoff
\[
g(S_T)=\1_{\{S_T>K\}}.
\]
\end{studentexercise}
\begin{shortsolution}
The payoff is one unit of cash exactly on the call exercise event. Therefore
\[
\boxed{
D(t,s)=e^{-r\tau}\Ncdf(d_-).}
\]
The discontinuity at maturity produces a sharply concentrated delta near
the strike as \(\tau\downarrow0\).
\end{shortsolution}

% ============================================================
\section{Greeks and the Replicating Strategy}
% ============================================================

\subsection{One identity simplifies every derivative}

For the call formula,
\[
C=s\Ncdf(d_+)-Ke^{-r\tau}\Ncdf(d_-).
\]
The normal densities satisfy
\[
\boxed{
s\phiN(d_+)=Ke^{-r\tau}\phiN(d_-).}
\]
Indeed, \(d_-=d_+-\sigma\sqrt{\tau}\) and direct simplification gives
\[
\frac{\phiN(d_-)}{\phiN(d_+)}
=\exp\left(\frac{d_+^2-d_-^2}{2}\right)
=\frac{se^{r\tau}}{K}.
\]

\subsection{Delta and gamma}

Because
\[
\partial_s d_+=\partial_s d_-
=\frac{1}{s\sigma\sqrt{\tau}},
\]
differentiating the call price gives
\[
\begin{aligned}
C_s
={}&\Ncdf(d_+)
+s\phiN(d_+)\partial_sd_+
-Ke^{-r\tau}\phiN(d_-)\partial_sd_-\\
={}&\Ncdf(d_+).
\end{aligned}
\]
Thus
\[
\boxed{\Delta_C=C_s=\Ncdf(d_+).}
\]
Differentiating again,
\[
\boxed{
\Gamma_C=C_{ss}
=\frac{\phiN(d_+)}{s\sigma\sqrt{\tau}}.}
\]

\begin{keyidea}{Hedging interpretation}
Delta is the number of shares in the replicating portfolio. Gamma measures
how quickly the delta changes. Large gamma makes continuous-time replication
especially sensitive to discrete rebalancing, transaction costs, jumps, and
volatility misspecification.
\end{keyidea}

\subsection{Vega, theta, and rho}

Differentiating with respect to volatility and using the density identity
gives
\[
\boxed{
\mathrm{Vega}_C=C_\sigma
=s\phiN(d_+)\sqrt{\tau}.}
\]
Vega is positive because the call payoff is convex.

Theta is most efficiently computed from the PDE:
\[
C_t=rC-rsC_s-\frac12\sigma^2s^2C_{ss}.
\]
Substitution yields
\[
\boxed{
\Theta_C=C_t
=-\frac{s\sigma\phiN(d_+)}{2\sqrt{\tau}}
-rKe^{-r\tau}\Ncdf(d_-).}
\]
Here theta means calendar-time sensitivity \(\partial C/\partial t\).

Finally,
\[
\partial_rd_+=\partial_rd_-=\frac{\sqrt{\tau}}{\sigma}.
\]
The density contributions cancel, giving
\[
\boxed{
\mathrm{Rho}_C=C_r
=K\tau e^{-r\tau}\Ncdf(d_-).}
\]

\begin{center}
\begin{tabular}{lcc}
\toprule
Greek & Call formula & Main role\\
\midrule
\(\Delta\) & \(\Ncdf(d_+)\) & stock hedge\\
\(\Gamma\) & \(\phiN(d_+)/(s\sigma\sqrt{\tau})\) & hedge curvature\\
Vega & \(s\phiN(d_+)\sqrt{\tau}\) & volatility sensitivity\\
\(\Theta\) &
\(-s\sigma\phiN(d_+)/(2\sqrt{\tau})-rKe^{-r\tau}\Ncdf(d_-)\)
& time decay\\
Rho & \(K\tau e^{-r\tau}\Ncdf(d_-)\) & rate sensitivity\\
\bottomrule
\end{tabular}
\end{center}

\begin{studentexercise}{put Greeks from parity}
Use
\[
P=C-s+Ke^{-r\tau}
\]
to derive the five put Greeks.
\end{studentexercise}
\begin{shortsolution}
\[
\begin{aligned}
\Delta_P&=\Ncdf(d_+)-1=-\Ncdf(-d_+),\\
\Gamma_P&=\Gamma_C,\qquad
\mathrm{Vega}_P=\mathrm{Vega}_C,\\
\Theta_P&=-\frac{s\sigma\phiN(d_+)}{2\sqrt{\tau}}
+rKe^{-r\tau}\Ncdf(-d_-),\\
\mathrm{Rho}_P&=-K\tau e^{-r\tau}\Ncdf(-d_-).
\end{aligned}
\]
\end{shortsolution}

\begin{warningbox}{Greeks are local model derivatives}
A Greek measures an infinitesimal sensitivity while holding the rest of the
model fixed. It is not a bound on finite losses. Near maturity or around a
payoff kink, higher-order terms and discrete trading can dominate a
first-order approximation.
\end{warningbox}

% ============================================================
\section{Discrete Hedging and Model Limitations}
% ============================================================

\subsection{The continuous hedge becomes a grid strategy}

On a grid \(0=t_0<\cdots<t_N=T\), a practical delta hedge uses
\[
\Delta_n=C_s(t_n,S_{t_n})
\]
over \((t_n,t_{n+1}]\). In discounted units,
\[
\bar S_t=e^{-rt}S_t,\qquad
\bar X_t=e^{-rt}X_t.
\]
The self-financing recursion is
\[
\boxed{
\bar X_{t_{n+1}}
=\bar X_{t_n}
+\Delta_n(\bar S_{t_{n+1}}-\bar S_{t_n}).}
\]
Starting from the Black--Scholes price,
\[
\boxed{
\bar X_T^N
=C(0,S_0)
+\sum_{n=0}^{N-1}
\Delta_n(\bar S_{t_{n+1}}-\bar S_{t_n}).}
\]
The terminal hedging error is
\[
\mathrm{PL}_T^N
=X_T^N-(S_T-K)^+.
\]
In the ideal model, the error converges to zero as the trading mesh tends to
zero under suitable conditions.

\begin{warningbox}{Why real hedging errors remain}
\begin{itemize}
  \item Trading occurs at discrete times and incurs costs.
  \item Volatility is not known and is not constant.
  \item Prices can jump and liquidity can disappear.
  \item Funding, collateral, dividends, and borrowing rates differ.
  \item The model may omit relevant state variables.
\end{itemize}
The Black--Scholes hedge is exact only inside its assumptions.
\end{warningbox}

\begin{studentexercise}{leading delta--gamma hedging error}
Over a short interval \(\Delta t\), freeze delta at the beginning of the
interval. Use a second-order Taylor expansion to identify the term that a
delta-only hedge does not remove.
\end{studentexercise}
\begin{shortsolution}
For a smooth option value,
\[
\Delta C
\approx C_t\Delta t+C_s\Delta S+\frac12C_{ss}(\Delta S)^2.
\]
After subtracting the frozen-delta stock gain \(C_s\Delta S\), the leading
unhedged price-move term is
\[
\boxed{\frac12\Gamma(\Delta S)^2.}
\]
The PDE balances its model-implied mean with theta and financing in
continuous time. Realized squared moves need not match the model-implied
variance over a discrete interval.
\end{shortsolution}

% ============================================================
\section{Guided Practice and Common Traps}
% ============================================================

\subsection{A complete replication drill}

\begin{studentexercise}{a smooth European payoff}
Let \(g(s)=s^2\).
\begin{enumerate}
  \item Use the general expectation formula to find \(p(t,s)\).
  \item Compute the replicating holdings \(\Delta_t,\beta_t\).
  \item Verify the PDE directly.
\end{enumerate}
\end{studentexercise}
\begin{shortsolution}
The power-payoff formula with \(m=2\) gives
\[
p(t,s)=s^2e^{(r+\sigma^2)(T-t)}.
\]
Hence
\[
\Delta_t=2S_te^{(r+\sigma^2)(T-t)}
\]
and
\[
\beta_t
=-\frac{S_t^2e^{(r+\sigma^2)(T-t)}}{B_t}.
\]
Direct substitution of
\[
p_t=-(r+\sigma^2)p,\qquad
p_s=2p/s,\qquad
p_{ss}=2p/s^2
\]
into the Black--Scholes PDE gives zero.
\end{shortsolution}

\subsection{Error diagnosis}

\begin{studentexercise}{correct four statements}
Correct the following claims:
\begin{enumerate}
  \item \(X_t=\Delta_tS_t+\beta_tB_t\) automatically implies
        self-financing.
  \item Rebalancing changes portfolio wealth.
  \item A delta-hedged claim always earns \(r\).
  \item The call exercise probability is \(\Ncdf(d_+)\).
\end{enumerate}
\end{studentexercise}
\begin{shortsolution}
\begin{enumerate}
  \item The value identity must be supplemented by
        \(\dd X=\Delta\dd S+\beta\dd B\).
  \item Self-financing rebalancing changes composition but not current wealth.
  \item The entire self-financing locally riskless portfolio earns \(r\)
        under no arbitrage; the option alone remains risky.
  \item Under the pricing dynamics, the exercise probability is
        \(\Ncdf(d_-)\).
\end{enumerate}
\end{shortsolution}

\subsection{End-of-lecture checklist}

\begin{keyidea}{Ten checks for a replication argument}
\begin{enumerate}
  \item State the stock and bank-account dynamics.
  \item Separate model assumptions from market assumptions.
  \item Specify the payoff and terminal condition.
  \item Define holdings and portfolio value.
  \item Impose predictability and self-financing.
  \item Apply \(\Ito\)'s formula to the candidate value.
  \item Match diffusion to identify delta.
  \item Match drift to derive the PDE.
  \item Check regularity, boundary behavior, and integrability.
  \item Verify that the strategy is self-financing and matches the payoff.
\end{enumerate}
\end{keyidea}

\begin{keyidea}{Final synthesis}
\begin{itemize}
  \item Self-financing is a cash-flow constraint, not merely a portfolio
        identity.
  \item Completeness allows Brownian risk to be replicated; no arbitrage makes
        the replication price unique.
  \item The physical drift cancels because the option and hedge carry the same
        stock exposure.
  \item Logarithmic coordinates turn the Black--Scholes PDE into a heat
        equation.
  \item The \(d_-\) term describes exercise under pricing dynamics; \(d_+\)
        comes from a stock-weighted Gaussian expectation.
  \item Delta is the hedge, gamma measures its instability, and the remaining
        Greeks diagnose parameter and carry sensitivity.
\end{itemize}
\end{keyidea}

\end{document}
